\section{Committed Inner Products}
\label{sec:inner-product}

The public linear functionals of \cref{sec:linear-functionals} use a public second factor.  We now pass to the first genuinely bilinear relation: both vectors are committed.  Let
\[
  a,b:\B^n\to\F
\]
be tables and let
\[
  S=\sum_{w\in\B^n}a(w)b(w)
\]
be their inner product.  By \cref{thm:middle-coefficient-inner-product},
\begin{equation}
  S=[X^{N-1}]V_a(X)V_b^{\star}(X).
  \label{eq:ip-middle-coefficient}
\end{equation}
The second factor is not public, but \cref{sec:virtual-reversal} shows that its Reed--Solomon word is already available virtually from the commitment to $b$:
\begin{equation}
  \ev_L(V_b^{\star})(\xi)
  =\xi^{N-1}\ev_L(V_b)(\xi^{-1}).
  \label{eq:ip-reversal-value}
\end{equation}
Thus the arbitrary committed inner product is an instance of the generic coefficient-extraction protocol with no additional commitment for the reversed second factor.

This section states the relation, gives the protocol explicitly, and derives its completeness and soundness from the general machinery of \cref{sec:generic-protocol}.  We also record the decoding consequence needed later when inner products are used as subrelations in lookup, sparse-matrix, and R1CS arguments.

\subsection{The inner-product relation}
\label{sec:ip-relation}

\begin{definition}[Hypercube inner product]
\label{def:hypercube-inner-product}
For tables $a,b:\B^n\to\F$, define
\begin{equation}
  \langle a,b\rangle_{\B^n}
  \defeq
  \sum_{w\in\B^n}a(w)b(w).
  \label{eq:hypercube-inner-product}
\end{equation}
When no confusion is possible we write simply $\langle a,b\rangle$.
\end{definition}

\begin{definition}[Committed-inner-product relation]
\label{def:ip-relation}
Fix $0<\delta\le(1-\rho)/2$.  The relation $\mathcal R^{\delta}_{\mathrm{IP}}$ consists of pairs
\[
  ((w_a,w_b;S),(a,b))
\]
with $w_a,w_b\in\F^L$, $S\in\F$, and tables $a,b:\B^n\to\F$ such that
\begin{align}
  \Delta_{\mathrm{fib}}(w_a,\Enc_T(a))&\le\delta,
  \label{eq:ip-relation-distance-a}\\
  \Delta_{\mathrm{fib}}(w_b,\Enc_T(b))&\le\delta,
  \label{eq:ip-relation-distance-b}\\
  S&=\langle a,b\rangle_{\B^n}.
  \label{eq:ip-relation-value}
\end{align}
The instance is $(w_a,w_b;S)$ and the witness is the pair $(a,b)$.
\end{definition}

The relation is meaningful because, in the unique-decoding radius, each source word determines at most one nearby table by \cref{prop:unique-table-near-word}.  In particular, once $w_a$ and $w_b$ are fixed, the value in \eqref{eq:ip-relation-value} is determined whenever both decodings exist.

\begin{lemma}[Inner products are generic coefficient claims]
\label{lem:ip-as-generic-ce}
Let $w_a,w_b\in\F^L$ and define the virtual reversal
\[
  w_b^{\star}(\xi)=\xi^{N-1}w_b(\xi^{-1}).
\]
Then $(w_a,w_b;S)$ has a witness in $\mathcal R^{\delta}_{\mathrm{IP}}$ if and only if $(w_a,w_b^{\star};S)$ has a witness in the generic coefficient-extraction relation $\mathcal R^{\delta}_{\mathrm{CE}}$ of \cref{def:generic-ce-relation}.
\end{lemma}

\begin{proof}
Suppose first that $(a,b)$ witnesses $\mathcal R^{\delta}_{\mathrm{IP}}$.  Put
\[
  U=V_a,
  \qquad
  K=V_b^{\star}.
\]
By \cref{lem:reversal-commutes-evaluation},
\[
  \ev_L(K)=\Enc_T(b)^{\star}.
\]
By the fibre-distance isometry \eqref{eq:reversal-fibre-distance},
\[
  \Delta_{\mathrm{fib}}(w_b^{\star},\ev_L(K))
  =\Delta_{\mathrm{fib}}(w_b,\Enc_T(b))
  \le\delta.
\]
The first distance condition is
\[
  \Delta_{\mathrm{fib}}(w_a,\ev_L(U))
  =\Delta_{\mathrm{fib}}(w_a,\Enc_T(a))
  \le\delta.
\]
Finally, \cref{thm:middle-coefficient-inner-product} gives
\[
  [X^{N-1}]UK
  =[X^{N-1}]V_aV_b^{\star}
  =\langle a,b\rangle_{\B^n}
  =S.
\]
Hence $(U,K)$ witnesses $\mathcal R^{\delta}_{\mathrm{CE}}$.

Conversely, suppose $(U,K)$ witnesses $\mathcal R^{\delta}_{\mathrm{CE}}$ for $(w_a,w_b^{\star};S)$.  By \cref{prop:table-isomorphism}, there is a unique table $a$ with $V_a=U$, and a unique table $b$ with
\[
  V_b=K^{\star};
\]
equivalently $V_b^{\star}=K$, since reversal is an involution by \cref{lem:reversal-properties}.  The first generic distance condition is exactly \eqref{eq:ip-relation-distance-a}.  For the second, use \cref{lem:reversal-commutes-evaluation} and \eqref{eq:reversal-fibre-distance}:
\[
  \Delta_{\mathrm{fib}}(w_b,\Enc_T(b))
  =\Delta_{\mathrm{fib}}(w_b^{\star},\ev_L(V_b^{\star}))
  =\Delta_{\mathrm{fib}}(w_b^{\star},\ev_L(K))
  \le\delta.
\]
The generic coefficient condition and \cref{thm:middle-coefficient-inner-product} give
\[
  S=[X^{N-1}]UK
   =[X^{N-1}]V_aV_b^{\star}
   =\langle a,b\rangle_{\B^n}.
\]
Thus $(a,b)$ witnesses $\mathcal R^{\delta}_{\mathrm{IP}}$.
\end{proof}

\subsection{The protocol}
\label{sec:ip-protocol}

For honest tables $a,b$, \cref{thm:universal-split} applied to $V_a$ and $V_b^{\star}$ gives unique $A,H\in\F[X]_{<N}$ satisfying
\begin{equation}
  XV_a(X)V_b^{\star}(X)
  =A(X)+SX^N+X^{N+1}H(X).
  \label{eq:ip-opening-identity}
\end{equation}
Only $A$ is sent.  The reversed second factor and $H$ are both virtual.

\begin{definition}[Inner-product protocol $\Pi_{\mathrm{IP}}$]
\label{def:ip-protocol}
The instance is $(w_a,w_b;S)$.
\begin{enumerate}[label=\textnormal{\arabic*.},leftmargin=*]
  \item The prover sends an oracle word $w_A\in\F^L$.  The verifier samples $\beta\in\F$.
  \item The verifier defines the virtual reversed word
  \begin{equation}
    w_b^{\star}(\xi)
    =\xi^{N-1}w_b(\xi^{-1})
    \qquad(\xi\in L),
    \label{eq:ip-virtual-reversal}
  \end{equation}
  and the virtual high-part word
  \begin{equation}
    h(\xi)
    =\frac{\xi w_a(\xi)w_b^{\star}(\xi)-w_A(\xi)-S\xi^N}{\xi^{N+1}}.
    \label{eq:ip-virtual-high}
  \end{equation}
  It forms
  \begin{equation}
    w_0^{(\beta)}
    =w_a+\beta w_b^{\star}+\beta^2w_A+\beta^3h.
    \label{eq:ip-batch}
  \end{equation}
  \item The parties run $\Pi_{\mathrm{FT}}[\ell,\kappa]$, or any sparse-level variant of \cref{def:committed-folding-levels}, on $w_0^{(\beta)}$.
\end{enumerate}
The verifier never receives either $w_b^{\star}$ or $h$ as a committed oracle.  At a query point $\xi$ it obtains $w_b^{\star}(\xi)$ by reading $w_b(\xi^{-1})$, and then computes $h(\xi)$ from \eqref{eq:ip-virtual-high}.
\end{definition}

\begin{remark}[Fibre-local query pattern]
\label{rem:ip-fibre-locality}
A folding query needs the batch at $\xi$ and $-\xi$.  The source values of $w_a$ and $w_A$ are read on the fibre $\{\xi,-\xi\}$, while the reversed source values are obtained from $w_b$ on the inverse fibre
\[
  \{\xi^{-1},-\xi^{-1}\}.
\]
By \cref{lem:inversion-preserves-fibres}, this is again one fibre of $L$.  Thus reversal changes the location of the second commitment query but does not destroy the fibre-local structure required by the folding backend.
\end{remark}

\begin{theorem}[Completeness of the inner-product protocol]
\label{thm:ip-completeness}
If
\[
  w_a=\Enc_T(a),
  \qquad
  w_b=\Enc_T(b),
  \qquad
  S=\langle a,b\rangle_{\B^n},
\]
then the honest prover is accepted by $\Pi_{\mathrm{IP}}$ with probability $1$.
\end{theorem}

\begin{proof}
By \eqref{eq:ip-middle-coefficient},
\[
  S=[X^{N-1}]V_aV_b^{\star}.
\]
Hence \cref{thm:universal-split} gives unique $A,H\in\F[X]_{<N}$ satisfying \eqref{eq:ip-opening-identity}, and the prover sends $w_A=\ev_L(A)$.  By \cref{lem:reversal-commutes-evaluation}, the virtual word \eqref{eq:ip-virtual-reversal} equals $\ev_L(V_b^{\star})$.  Evaluating \eqref{eq:ip-opening-identity} at $\xi\in L$ gives
\[
  H(\xi)
  =\frac{\xi V_a(\xi)V_b^{\star}(\xi)-A(\xi)-S\xi^N}{\xi^{N+1}},
\]
so \eqref{eq:ip-virtual-high} equals $\ev_L(H)$.  Therefore, for every $\beta\in\F$,
\[
  w_0^{(\beta)}
  =\ev_L\!\left(V_a+\beta V_b^{\star}+\beta^2A+\beta^3H\right)
  \in C_0.
\]
Completeness follows from \cref{prop:folding-completeness}.
\end{proof}

\subsection{Soundness}
\label{sec:ip-soundness}

The protocol is exactly the generic coefficient-extraction protocol applied to the pair $(w_a,w_b^{\star})$.  The only point that must be transported back to the original instance is the decoding of the second commitment; the reversal isometry of \cref{lem:inversion-preserves-fibres} does precisely this.

\begin{theorem}[Soundness of committed inner products]
\label{thm:ip-soundness}
Assume
\begin{equation}
  2N<(1-\delta)M.
  \label{eq:ip-degree-margin}
\end{equation}
If $(w_a,w_b;S)$ has no witness in $\mathcal R^{\delta}_{\mathrm{IP}}$, then every prover is accepted by $\Pi_{\mathrm{IP}}$ with probability at most
\begin{equation}
  \varepsilon_{\mathrm{IP}}
  \defeq
  \frac{3M}{|\F|}
  +\varepsilon_{\mathrm{fold}}
  +(1-\delta)^\kappa.
  \label{eq:ip-soundness-bound}
\end{equation}
The same bound holds for every sparse-level folding variant covered by \cref{prop:virtual-levels}.
\end{theorem}

\begin{proof}
Suppose the original instance has no witness in $\mathcal R^{\delta}_{\mathrm{IP}}$.  By \cref{lem:ip-as-generic-ce}, the transformed instance
\[
  (w_a,w_b^{\star};S)
\]
has no witness in $\mathcal R^{\delta}_{\mathrm{CE}}$.  The protocol of \cref{def:ip-protocol} is exactly $\Pi_{\mathrm{CE}}$ of \cref{def:generic-ce-protocol} on this transformed instance: the second source word is evaluated virtually using \eqref{eq:ip-virtual-reversal}, while the generic high-part word is exactly \eqref{eq:ip-virtual-high}.  Therefore \cref{thm:generic-ce-soundness} gives \eqref{eq:ip-soundness-bound}.  The sparse-level statement follows from the same theorem.
\end{proof}

\begin{remark}[Why the field term is $3M/|\F|$]
\label{rem:ip-three-M}
Unlike a public linear functional, the second factor is itself committed and must be protected by correlated agreement.  The batch \eqref{eq:ip-batch} contains four degree-$<N$ words,
\[
  w_a,\quad w_b^{\star},\quad w_A,\quad h,
\]
and is a degree-$3$ curve in $\beta$.  Thus the correlated-agreement step uses $t=3$, which is the origin of the term $3M/|\F|$.  The reversal is virtual, but cryptographically it remains an independently decoded source word.
\end{remark}

\subsection{Decoding, uniqueness, and binding consequences}
\label{sec:ip-decoding}

The soundness theorem is stated as a relation statement.  For later composition it is useful to record explicitly what a witness means when the two commitment words are individually decodable.

\begin{proposition}[Decoded value is unique]
\label{prop:ip-decoded-value-unique}
Assume $0<\delta\le(1-\rho)/2$.  Let $w_a,w_b\in\F^L$.  There is at most one pair of tables $(a,b)$ satisfying
\[
  \Delta_{\mathrm{fib}}(w_a,\Enc_T(a))\le\delta,
  \qquad
  \Delta_{\mathrm{fib}}(w_b,\Enc_T(b))\le\delta.
\]
Consequently, there is at most one $S\in\F$ for which $(w_a,w_b;S)$ has a witness in $\mathcal R^{\delta}_{\mathrm{IP}}$.
\end{proposition}

\begin{proof}
Uniqueness of $a$ and $b$ follows independently from \cref{prop:unique-table-near-word}.  Once $a$ and $b$ are fixed, \eqref{eq:ip-relation-value} forces
\[
  S=\langle a,b\rangle_{\B^n},
\]
so $S$ is unique.
\end{proof}

\begin{corollary}[Evaluation binding for an inner-product value]
\label{cor:ip-binding}
Under the assumptions of \cref{thm:ip-soundness}, fix $w_a,w_b$ and two distinct values $S\ne S'$.  It is impossible that both instances $(w_a,w_b;S)$ and $(w_a,w_b;S')$ have witnesses in $\mathcal R^{\delta}_{\mathrm{IP}}$.  Moreover, if one of the two instances is false, every prover for that instance is accepted with probability at most $\varepsilon_{\mathrm{IP}}$.
\end{corollary}

\begin{proof}
The first statement is \cref{prop:ip-decoded-value-unique}.  The second is \cref{thm:ip-soundness}.
\end{proof}

\begin{proposition}[Polynomial-time witness extraction from decoded source words]
\label{prop:ip-decoding-extractor}
Assume $0<\delta\le(1-\rho)/2$.  Suppose a polynomial-time Reed--Solomon unique decoder is available for radius $\delta$.  Given $w_a,w_b$, one can in time polynomial in $M$ either reject because one source word has no codeword within distance $\delta$, or recover the unique tables $a,b$ satisfying the two distance conditions of $\mathcal R^{\delta}_{\mathrm{IP}}$.  The claimed value is then checked by computing $\langle a,b\rangle_{\B^n}$.
\end{proposition}

\begin{proof}
Apply the Reed--Solomon decoder separately to $w_a$ and $w_b$.  If it returns codewords $c_a,c_b\in C_0$, recover their unique degree-$<N$ polynomials using the inverse evaluation map of \cref{prop:rs-parameters}.  By \cref{prop:table-isomorphism}, their coefficients define unique tables $a,b$.  Check the two fibre distances explicitly and then compute the sum in \eqref{eq:hypercube-inner-product}.  All operations require time polynomial in $M$.
\end{proof}

The proposition is the algebraic extractor used later.  The round-by-round knowledge state needed for non-interactive compilation will be developed globally in the later soundness section; no additional algebraic extraction argument is required for inner products.

\subsection{Costs and relation to degree-two Sumcheck}
\label{sec:ip-costs}

The identity \eqref{eq:ip-middle-coefficient} replaces the standard degree-two Boolean-hypercube sum
\begin{equation}
  S=\sum_{w\in\B^n}a(w)b(w)
  \label{eq:ip-degree-two-sum}
\end{equation}
by one coefficient-extraction relation.  A conventional Sumcheck proof of \eqref{eq:ip-degree-two-sum} recursively eliminates the $n$ Boolean variables and terminates in evaluations of the multilinear extensions of $a$ and $b$ at one random point.  Here there are no Sumcheck round polynomials.  Instead, the prover computes the product $V_aV_b^{\star}$, splits it around degree $N$, commits to the low part $A$, and proves one Reed--Solomon proximity statement.

\begin{proposition}[Local verifier work]
\label{prop:ip-local-verifier-work}
Apart from the folding backend, one queried fibre for $\Pi_{\mathrm{IP}}$ requires:
\begin{enumerate}[label=(\roman*),leftmargin=*]
  \item the values of $w_a$ and $w_A$ on one fibre of $L$;
  \item the values of $w_b$ on the corresponding inverse fibre;
  \item $O(1)$ field operations per queried point to compute the reversal \eqref{eq:ip-virtual-reversal}, the high-part value \eqref{eq:ip-virtual-high}, and the batch \eqref{eq:ip-batch}.
\end{enumerate}
In particular, unlike multilinear evaluation in \cref{sec:table-mle-opening}, no $O(n)$ product-form kernel evaluation is needed.
\end{proposition}

\begin{proof}
Items (i) and (ii) follow from \cref{rem:ip-fibre-locality}.  Given the source values, \eqref{eq:ip-virtual-reversal} uses one multiplication by the known scalar $\xi^{N-1}$, and \eqref{eq:ip-virtual-high} and \eqref{eq:ip-batch} use a constant number of field operations.  Powers of the queried point are determined by the domain point and may be precomputed or updated along the folding query path.  No evaluation of a length-$n$ structured kernel appears.
\end{proof}

\begin{remark}[Prover arithmetic]
\label{rem:ip-prover-arithmetic}
The additional first-round algebra beyond the table commitments is the product $V_aV_b^{\star}$ and the evaluation/commitment of its low part $A$.  With FFT/NTT multiplication, the product is quasi-linear in $N$ over a suitable smooth domain.  Exact implementation costs are deliberately postponed to the benchmark section: the mathematical comparison with Sumcheck depends on whether product computation, low-part encoding, Merkle commitments, and the common folding backend are fused or shared with surrounding relations.
\end{remark}

\begin{remark}[Why this primitive matters later]
\label{rem:ip-later-applications}
The inner-product protocol is the committed analogue of the public-linear-functional protocol.  Later sections use it as a component rather than as an isolated end goal.  In particular:
\begin{itemize}[leftmargin=*]
  \item degree-two Sumcheck relations of the form $\sum_w a(w)b(w)$ are immediate instances;
  \item logarithmic-derivative lookup and multiset arguments reduce their global consistency equation to inner products;
  \item sparse matrix--vector products reduce, after random row compression, to an inner product between a compressed matrix row and the committed vector;
  \item bounded-degree pointwise arithmetic expressions are reduced to Hadamard relations followed by a final inner product or public linear functional.
\end{itemize}
Thus \cref{sec:inner-product} is the first step from coefficient extraction as an opening technique to coefficient extraction as a general proof-system primitive.
\end{remark}
